% =====================================================================
%  Methods section -- Ghost Tasking for physics-consistent Gaussian
%  processes.  Drop-in LaTeX; requires amsmath, amssymb, bm.
%  Figure placeholders are described in comments -- no figures included.
% =====================================================================

\section{Methods}
\label{sec:methods}

\subsection{General framework}
\label{sec:framework}

We model each problem with a multi-output Gaussian process (GP) whose
covariance is constructed so that every sample path satisfies the
governing linear (system of) differential equation(s) exactly.  Let
$L$ denote a matrix of linear differential operators with
(possibly unknown) constant or polynomial coefficients, acting on a
vector of latent fields, and write the homogeneous constraint as
$L\,\bm{u}=\bm{0}$.  Following the parametrization approach to
physics-consistent GPs (PCGP)~\cite{XXX_PCGP}, we seek a matrix of
operators $B$ that \emph{parametrizes} the solution space, i.e.
\begin{equation}
  L\,B \;=\; \bm{0},
  \qquad
  \bm{u} \;=\; B\,g ,
  \label{eq:parametrization}
\end{equation}
so that for \emph{any} scalar (or low-dimensional) latent function $g$
the field $\bm{u}=Bg$ is an exact solution.  Placing a GP prior on $g$
with base kernel $k_\theta$ (a squared-exponential kernel with
amplitude and length scale $\theta$ throughout) induces a matrix-valued
covariance on the tasks,
\begin{equation}
  K(\bm{x},\bm{x}')
  \;=\;
  B_{\bm{x}}\,k_\theta(\bm{x},\bm{x}')\,B_{\bm{x}'}^{\!\top},
  \label{eq:matrix-kernel}
\end{equation}
where $B_{\bm{x}}$ applies the operators in $B$ in the argument $\bm{x}$.
The parametrization $B$ is obtained as a syzygy of $L$ in the Weyl
algebra and is computed symbolically once per problem in
\textsc{Macaulay2}~\cite{XXX_M2}; the resulting closed-form entries of
$K$ are exported to a \textsc{GpyTorch} kernel.  Boundary/initial
conditions, source terms and field observations enter as additional
\emph{tasks} (rows of $B$), evaluated at their respective input
locations.

\paragraph{Ghost Tasking (GT).}
The contribution of this work is to extend the parametrization
\eqref{eq:parametrization} from homogeneous problems to inhomogeneous
ones, and to gauge/constraint handling, by augmenting the operator
$L$ with additional columns---\emph{ghost tasks}---before computing the
syzygy.  Concretely we form $\tilde L=[\,L\mid C\,]$, where the columns
$C$ encode the inhomogeneity (a source term) or an auxiliary constraint
(e.g.\ a gauge), and compute the parametrization $\tilde B$ of
$\tilde L$.  The ghost tasks let the latent process generate both the
solution and its source/constraint \emph{consistently}, so the
inhomogeneous relation holds exactly for every $g$ rather than being
imposed only at the data.  This (i) reduces the number of independent
latent processes, (ii) embeds the source/gauge in the kernel, and
(iii) improves the conditioning of $K$ (Sec.~\ref{sec:ex3}).

\paragraph{Comparison: physics-informed GP (PIGP).}
We benchmark GT against the physics-informed parametrization used in
prior work~\cite{XXX_PIGP}, in which the latent process \emph{is} the
solution (the leading block of $B$ is the identity) and the operator
image $L\bm{u}$ is appended as extra tasks whose targets are the
observed source terms.  In PIGP the physics is therefore enforced
\emph{through the data} (residual tasks driven to the source values),
whereas in GT it is built into the kernel structure.  Both methods use
the identical base kernel, training procedure and data, so differences
in accuracy, parameter recovery and conditioning are attributable to
the parametrization alone.

\paragraph{Training and inference.}
Given data $\bm{y}$ at inputs $X$ with a fixed-variance Gaussian
likelihood (homoscedastic noise $\sigma^2$, exact constraints assigned a
negligible jitter), the kernel and physical parameters are obtained by
maximizing the exact marginal log-likelihood with the Adam optimizer
(learning rate $0.1$, $500$ iterations).  Forward problems predict the
fields and derived quantities at test locations from initial/boundary
data and known sources; inverse problems additionally treat the
physical parameter(s) as trainable and recover them from noisy
observations.  Posterior uncertainty on the recovered parameters is
quantified by a Laplace approximation about the maximum-a-posteriori
estimate (the Hessian of the negative log-likelihood).  As a
ground-truth reference for the inverse problems we evaluate the exact
(brute-force) marginal likelihood on a dense grid of the physical
parameter(s), against which both the point estimates and the Laplace
covariance are validated.  Accuracy is reported as the (R)MSE between
the posterior mean and the analytic solution on a held-out test grid,
together with the recovered parameter values, their Laplace
uncertainties and the wall-clock training time.

% ---------------------------------------------------------------------
\subsection{Experiment 1: pedagogical coupled transport system}
\label{sec:ex1}

As a controlled, low-dimensional illustration we consider a coupled
first-order linear system for $(u,v)$ on $(x,t)\in[0,1]^2$,
\begin{equation}
  \bigl(x\,\partial_x+\partial_t\bigr)u + a\,\partial_x v = 0,
  \qquad
  \bigl(x\,\partial_x+\partial_t\bigr)v = 0,
  \label{eq:ex1-pde}
\end{equation}
with operator matrix
$L=\bigl[\begin{smallmatrix} x\partial_x+\partial_t & a\,\partial_x\\
0 & x\partial_x+\partial_t\end{smallmatrix}\bigr]$ and scalar coupling
parameter $a$.  The system admits the closed-form solution
$v(x,t)=v_0(x e^{-t})$,
$u(x,t)=u_0(x e^{-t})-a\,(1-e^{-t})\,v_0'(x e^{-t})$,
which we use to generate data and to evaluate errors; here
$u_0\equiv0$, $v_0(x)=3\sin(2x)$ and the true value is $a=2$.

The GT parametrization is the single-column syzygy of $L$ (one latent
process, three tasks $u,v$ and the homogeneous constraint); the ghost
variant augments $L$ with a source column to handle the inhomogeneous
case.  The PIGP parametrization uses two latent processes and appends
the two PDE residuals as tasks (four tasks total).  Training data lie on
$n\times n$ grids with $n\in\{10,15\}$; the inverse problem is solved at
noise levels $\sigma\in\{0.001,0.01,0.1,0.2\}$, recovering $a$.  Because
the problem is two-dimensional we can compute the exact posterior over
$a$ on a fine grid, providing an unambiguous reference for the
GT/PIGP point estimates and Laplace uncertainties.

\paragraph{Suggested figures.}
\begin{itemize}
  \item \textbf{Fig.~1 (forward solution).} Side-by-side surface/contour
        plots of the GT and PIGP posterior mean for $u$ and $v$ over
        $(x,t)$, with the analytic solution and the pointwise error;
        overlay training points. Demonstrates exact constraint
        satisfaction and the effect of the extrapolation point.
  \item \textbf{Fig.~2 (parameter recovery).} The exact 1-D posterior
        over $a$ (grid evaluation) overlaid with the GT and PIGP MAP
        estimate $\pm$ Laplace standard deviation, for the four noise
        levels (small multiples). Shows calibration of the Laplace
        approximation.
  \item \textbf{Fig.~3 (accuracy vs.\ noise).} Test RMSE and
        $|\hat a-a|$ as a function of $\sigma$ for both methods and both
        grid sizes (line plot with markers).
\end{itemize}

% ---------------------------------------------------------------------
\subsection{Experiment 2: forced triple pendulum (boundary-value problem)}
\label{sec:ex2}

We next consider three identical linear pendula of length $\ell$ under
gravity $g$, driven by a common but unknown forcing $u(t)$,
\begin{equation}
  \ell\,\ddot\theta_i(t) + g\,\theta_i(t) = u(t),
  \qquad i=1,2,3,
  \label{eq:ex2-ode}
\end{equation}
on $t\in[0,t_1]$ with two-point boundary conditions
$\theta_i(0),\theta_i(t_1)$ prescribed.  The operator acts on
$(\theta_1,\theta_2,\theta_3,u)$ as
$L=\bigl[(\ell\partial_t^2+g)\,I_3 \mid -\bm{1}\bigr]$, coupling the
three trajectories through the shared drive $u$.  The true parameters
are $g=9.81$ and $\ell=2.5$; data are generated from the analytic
modal solution of \eqref{eq:ex2-ode} (sinusoidal drive
$u(t)=A\sin(\omega_u t)$).

The GT parametrization (one latent process; tasks
$\theta_1,\theta_2,\theta_3,u$ plus a ghost task) augments $L$ with a
forcing column $[\,\ell,\ \ell\partial_t,\ \ell t\,]^{\top}$, so the
common drive is generated consistently with the trajectories.  PIGP uses
three latent processes and appends the three equation residuals as
tasks.  Only sparse boundary observations of $\theta_i$ (at $t=0,t_1$)
and samples of the drive are provided.  We study sample sizes
$n\in\{3,5,10,25\}$ and noise levels
$\sigma\in\{0,0.01,0.1,0.2\}$, and jointly recover $(g,\ell)$.

\paragraph{Suggested figures.}
\begin{itemize}
  \item \textbf{Fig.~4 (reconstructed dynamics).} The three trajectories
        $\theta_1,\theta_2,\theta_3$ and the inferred forcing $u(t)$ over
        $[0,t_1]$: GT and PIGP posterior mean with credible band,
        analytic solution, boundary/observation points. Highlights
        reconstruction of the latent drive from boundary data.
  \item \textbf{Fig.~5 (joint parameter posterior).} 2-D contour of the
        exact $(g,\ell)$ likelihood from the grid evaluation, with GT and
        PIGP MAP estimates and Laplace covariance ellipses overlaid.
  \item \textbf{Fig.~6 (data-efficiency).} Test RMSE and parameter error
        vs.\ $n$ (and vs.\ $\sigma$) for GT and PIGP, illustrating
        robustness in the data-scarce regime.
\end{itemize}

% ---------------------------------------------------------------------
\subsection{Experiment 3: 3-D magnetostatics with spatially varying
            reluctivity}
\label{sec:ex3}

The final, most demanding example is 3-D magnetostatics in the
vector-potential formulation,
\begin{equation}
  \nabla\times\!\bigl(\nu\,\nabla\times\bm{A}\bigr) = \bm{J},
  \qquad
  \bm{B}=\nabla\times\bm{A},
  \label{eq:ex3-pde}
\end{equation}
on $\Omega=[0,1]^2\times[1,2]$ with a spatially varying, anisotropic
reluctivity tensor $\nu=\mathrm{diag}(\nu_0,\nu_0,z)$ and Dirichlet
conditions $\bm{A}=\bm{0}$ on the domain faces.  The operator
$L=\mathrm{curl}\,\nu\,\mathrm{curl}$ acts on $\bm{A}=(A_1,A_2,A_3)$;
because $L$ has a non-trivial null space (gradient fields), a gauge
condition must be supplied.  A manufactured analytic solution with the
corresponding source $\bm{J}$ is used as ground truth, and the unknown
scalar reluctivity $\nu_0$ ($\nu_0=1$) is recovered in the inverse
setting.

The GT parametrization augments $L$ with a source column and a gauge
\emph{ghost task}, yielding a single-latent-process kernel whose
sample paths satisfy both \eqref{eq:ex3-pde} and the gauge exactly; the
tasks are the three components of $\bm{A}$, the source $J_x$, the gauge
(target $0$) and the three components of $\bm{B}$.  PIGP instead models
$\bm{A}$ directly and appends the three components of $L\bm{A}$ as
residual tasks driven to $(J_x,0,0)$, together with the curl tasks for
$\bm{B}$.  Forward runs use boundary points on the faces plus the source
on an $n^3$ grid ($n=6$); inverse runs recover $\nu_0$ from noisy
field and source measurements over $n\in\{3,5,6,10\}$ and
$\sigma\in\{0.001,0.01,0.1,0.2\}$, with several random restarts.  This
example also exposes the conditioning advantage of GT: because the gauge
and source are absorbed into the kernel, the spectrum of $K$ decays more
slowly than for PIGP, which we report directly.

\paragraph{Suggested figures.}
\begin{itemize}
  \item \textbf{Fig.~7 (field reconstruction).} Slices (or quiver/stream
        plots) of the predicted $\bm{B}$ field on representative planes
        in $\Omega$ for GT and PIGP, with the analytic field and the
        pointwise error. Shows divergence-free, gauge-consistent
        reconstruction.
  \item \textbf{Fig.~8 (conditioning).} Sorted eigenvalue spectrum of the
        training covariance $K$ (log scale) for GT vs.\ PIGP at fixed
        grid/noise. Central evidence that the ghost-task kernel is
        better conditioned.
  \item \textbf{Fig.~9 (inverse $\nu_0$).} Recovered $\nu_0$ with Laplace
        uncertainty vs.\ noise level (and vs.\ $n$), GT vs.\ PIGP,
        against the exact grid posterior; include a box/violin over
        random restarts.
  \item \textbf{Fig.~10 (cost vs.\ accuracy).} Test RMSE vs.\ training
        time (or vs.\ number of tasks/latent processes) for GT and PIGP,
        summarizing the efficiency gain.
\end{itemize}

% ---------------------------------------------------------------------
\subsection{Implementation}
All models are implemented in \textsc{GpyTorch}/\textsc{PyTorch} in
double precision; symbolic parametrizations are derived with
\textsc{SymPy} and the syzygies with \textsc{Macaulay2}.  Code and the
data-generation scripts for all three experiments are available
in the accompanying repository~\cite{XXX_repo}.
